% Шаблон документа для курса "Системы научных публикаций"
% Семинар 10. LaTeX. Структура документа. Базовые элементы форматирования
\documentclass[a4paper,14 pt]{article}

%%% Работа с русским языком
\usepackage{cmap}          % Поиск в PDF
\usepackage{mathtext}         % Русские буквы в фомулах
\usepackage[T2A]{fontenc}      % Кодировка
\usepackage[utf8]{inputenc}      % Кодировка исходного текста
\usepackage[english,russian]{babel}  % Локализация и переносы
\usepackage{indentfirst}            % Для красной строки первого абзаца

\usepackage{extsizes}               % Возможность сделать 14-й шрифт
\usepackage{geometry}               % Простой способ задавать поля
\geometry{top=25mm}
\geometry{bottom=25mm}
\geometry{left=20mm}
\geometry{right=20mm}

\usepackage{fancyhdr}                   % Колонтитулы
\pagestyle{fancy}
\renewcommand{\headrulewidth}{0.2mm}  % Толщина линейки, отчеркивающей
% \lfoot{Нижний левый}                % верхний колонтитул
% \rfoot{Нижний правый}
\rhead{}
\chead[C]{МГТУ им. Н.Э. Баумана}
\lhead{}
\cfoot{Страница \thepage\ из \pageref{LastPage}} % По умолчанию здесь номер страницы (\thepage)


\usepackage{setspace} % Межстрочное расстояние
\onehalfspacing      % Межстрочное расстояние 1.5
%\doublespacing       % Межстрочное расстояние 2
%\singlespacing       % Межстрочное расстояние 1

\usepackage{lastpage} % Узнать, сколько всего страниц в документе

\usepackage{soul}     % Модификаторы начертания

\usepackage{multicol} % Несколько колонок

\usepackage{lipsum}   % Генерация случайного текста

%\renewcommand{\familydefault}{\sfdefault} % Начертание шрифта без засечек

\usepackage{hyperref}
\usepackage[usenames,dvipsnames,svgnames,table,rgb]{xcolor}
\hypersetup{        % Гиперссылки
  unicode=true,           % Русские буквы в раздела PDF
  pdftitle={Заголовок},   % Заголовок
  pdfauthor={Автор},      % Автор
  pdfsubject={Тема},      % Тема
  pdfcreator={Создатель}, % Создатель
  pdfproducer={Производитель}, % Производитель
  pdfkeywords={keyword1} {key2} {key3}, % Ключевые слова
  colorlinks=true,         % false: ссылки в рамках; true: цветные ссылки
  linkcolor=black,        % Внутренние ссылки
  citecolor=green,        % Ссылки на библиографию
  filecolor=magenta,      % Ссылки на файлы
  urlcolor=red           % Ссылки на URL
}

\usepackage{lipsum} % Генерация случайного текста

%%% Заголовок
\author{ФИО}
\title{\LaTeX. Структура документа. Базовые элементы форматирования}
\date{\today}

\begin{document} % Конец преамбулы, начало документа

\begin{titlepage}
  \thispagestyle{empty}

  \begin{center}
    {\small
      Министерство науки и высшего образования Российской Федерации\\
      Московский государственный технический университет\\
      имени Н.Э. Баумана\\[0.4cm]
    }

    {\normalsize
      Факультет ИУ «Информатика и системы управления»\\
      Кафедра ИУ-1 «Системы автоматического управления»
    }

    \vfill

    {\Large\bfseries ОТЧЁТ}\\[0.4cm]
    {\large по семинару № 10}\\[0.3cm]
    {\large по дисциплине}\\[0.2cm]
    {\Large\bfseries «Системы научных публикаций»}

    \vfill
  \end{center}

  \begin{flushright}
    Соин Андрей Дмитриевич\\
    Величко  Фёдор Тимофеевич\\
    Группа ИУ1-51Б\\[0.2cm]

  \end{flushright}

  \vfill

  \begin{center}

    2026
  \end{center}
\end{titlepage}

\tableofcontents

\section{Размер и начертание текста}

\label{sec:text}

\subsubsection{Large текст}

\Large Очень большой текст.

\subsubsection{Huge текст}

\begin{huge}
  Огромный текст.
\end{huge}

\subsubsection{Small подподраздел}

\small Маленький текст

\subsubsection{tiny подраздел}

\begin{tiny}
  Очень маленький текст
\end{tiny}

\section{Второй раздел}

Ну тут короче всякие виды выделения текста. \textbf{ Жирный текст.} \textit{Курсив. } \underline{Подчеркнутый. } \st{Зачеркнутый. } \emph{Выделение. }

\subsection[Подраздельчик]{Сокращенный подраздел}

\subsection*{Ненумерованый подраздел}

\section{Третий раздел}

\section{Списки}

\subsection{Маркированный список}

Маркированный список используется для перечисления элементов без нумерации:

\begin{itemize}
  \item Первый пункт списка.
  \item Второй пункт списка.
  \item[-] Третий пункт с изменённым маркером.
  \item Четвёртый пункт списка.
\end{itemize}

\subsection{Нумерованный список}

Нумерованный список используется для перечисления элементов по порядку:

\begin{enumerate}
  \item Первый пункт.
  \item Второй пункт.
  \item Третий пункт.
  \item Четвёртый пункт.
\end{enumerate}

\subsection{Вложенный список из трёх уровней}

Пример вложенного списка, содержащего три уровня:

\begin{itemize}
  \item Первый уровень — маркированный список.

    \begin{itemize}
      \item Второй уровень — маркированный список.
      \item Второй пункт второго уровня.

        \begin{enumerate}
          \item Третий уровень — нумерованный список.
          \item Второй пункт третьего уровня.
          \item Третий пункт третьего уровня.
        \end{enumerate}

      \item Ещё один пункт второго уровня.
    \end{itemize}

  \item Ещё один пункт первого уровня.
\end{itemize}




\section{Колонки}

Для размещения текста в нескольких колонках
в LaTeX используется окружение multicols.
Оно позволяет разделить текст на несколько
вертикальных колонок.

\begin{multicols}{3}

  LaTeX является системой компьютерной вёрстки,
  предназначенной для создания научных статей,
  отчётов, книг и технической документации.

  Одним из главных преимуществ LaTeX является
  возможность автоматического форматирования
  документа. Пользователь задаёт структуру,
  а система самостоятельно выполняет вёрстку.

  LaTeX широко применяется в научной среде,
  поскольку позволяет создавать математические
  формулы, таблицы, списки и библиографию.

  Использование нескольких колонок позволяет
  эффективно размещать текст на странице
  и делает документ более компактным.

\end{multicols}



\section{Ссылки}

В LaTeX можно создавать ссылки на другие
разделы документа и внешние интернет-ресурсы.

\subsection{Ссылка на раздел документа}

Раздел «Размер и начертание текста»
имеет номер \ref{sec:text} и находится
на странице \pageref{sec:text}.

Таким образом, с помощью команд ref и pageref
можно автоматически получать номер раздела
и страницу, на которой он расположен.

\subsection{Ссылка на сайт МГТУ им. Н.Э. Баумана}

Ссылка с помощью команды url:

\url{https://bmstu.ru}

Ссылка с помощью команды href:

\href{https://bmstu.ru}
{Официальный сайт МГТУ им. Н.Э. Баумана}.


\end{document}